% 解説プリントのコメント（saiten が書く。./tools/kaisetsu.ps1 がこのファイルからプリントを作る）
%% matome
\begin{daiji}
  \item 「連続量として扱える理由」は，準位間隔 $\Delta E$ が $\kB T$ より桁違いに小さいことを，粗視化の条件 $\Delta E_{\mathrm{level}}\ll\delta U\ll U$ を満たす幅 $\delta U$ が選べることに結びつけて初めて説明になる（問題1 (3)）。
  \item 状態の数え上げでは，答えの準位だけでなく「その間に別の値が無い」ことまで示して，漏れなく調べたことを伝える（問題1 (1) の和10）。
  \item $W\propto U^{N/2}$ で $N$ が巨大なとき，$U$ 以下の状態のほとんどはエネルギーの上端のごく薄い殻にある。これが $W$ と $w=\Omega\,\delta U$ が同じエントロピーを与える理由で，問題3 (5) はここを確かめる問いである（未着手）。
\end{daiji}
\begin{maemo}
  \item 理由の書き落とし：統計力学1-1（旧版）問題2 (4)（2026-10-27）で「$2^{N}$ になる理由」，統計力学1-1 問題3 (4)（2026-11-11）で $\Delta U=0$ の理由が無かった。今回も問題1 (3) で，結論「連続的に扱える」と根拠「間隔が小さい」の間をつなぐ一言（$\delta U$ が選べる）が抜けた。「理由を述べよ」の問いでは，要点整理のどの条件・定義に当てはまるのかを名指しして書く。
\end{maemo}
%% mondai 1
\begin{kotaeran}
  \machigai{1}{○}{記述} $(1,1,1)$ で $E/E_{0}=3$（縮退度1），$(2,1,1)$ の並べ替えで6（3），$(2,2,1)$ で9（3），$(3,1,1)$ で11（3）と，値も縮退度も正しい。ただし，9 の次が 11 であること，すなわち和10 が $1,4,9$ から重複を許して3つ選んだ和では作れないこと（次の候補 $(2,2,2)$ は12）が式にも文にも無い。並べた4つが「低い方から4つ」であることの根拠が欠けている。
  \machigai{3}{○}{記述} $\Delta E/(\kB T)=\num{6.0e-17}$ は正しい。理由は「$\Delta E\ll\kB T$ なので，温度に対してエネルギー変化の間隔が十分小さいため連続的に扱える」と書いた。間隔が小さいことと「連続量として扱える」の間に，粗視化の条件 $\Delta E_{\mathrm{level}}\ll\delta U\ll U$ を満たす幅 $\delta U$ が選べる（幅 $\delta U$ の中に非常に多くの準位が入り，離散性が見えない）という一言が要る。
\end{kotaeran}
\begin{rikai}
  \item ミクロカノニカル分配関数 $w$ は幅 $\delta U$ の中の固有状態の数である。$\Delta E_{\mathrm{level}}\ll\delta U$ なら幅の中に多数の準位が入り，$w$ は $\delta U$ に比例する滑らかな量になる。$\delta U\ll U$ なら $U$ はほぼ一点に決まる。この両方が同時に満たせることが「エネルギーを連続量として扱える」の中身で，$\Delta E/(\kB T)=\num{6.0e-17}$ はその両立が楽にできることを示している。
  \item 数え上げの確かめ方：和の候補を小さい順に $3,4,5,\dots$ と一つずつ当たり，作れない値（ここでは 4,5,7,8,10）を飛ばしたことを一言書く。各 $n_{i}^{2}$ は $1,4,9,\dots$ なので，3つの和として作れるかどうかを値ごとに具体的に確かめる（10 は $1+9$ に残り $0$ が要り，$4+4+$ に残り $2$ が要るので作れない）。
\end{rikai}
%% mondai 3
\begin{kotaeran}
  \machigai{5}{×}{未着手} (4) までで終わっていて，(5) の答案が無い。
\end{kotaeran}
\begin{rikai}
  \item (4) で $\Omega=(N/(2U))\,W$ と求めたことが (5) の土台になる。$W\propto U^{N/2}$ なので，比 $W(U(1-\eta))/W(U)$ は $U$ によらず $\eta$ と $N$ だけで決まる。$N\sim10^{23}$ だと，$\eta$ がごく小さくても指数 $N/2$ のせいで比は桁違いに小さくなる。
  \item その結果「$U$ 以下の状態のほとんどは上端の薄い殻 $U(1-\eta)<H\le U$ にある」と言えるので，殻の中だけを数える $w=\Omega\,\delta U$ と全体を数える $W$ とで，エントロピーの $N$ に比例する部分が変わらない。(4) の差 $\kB\log\{N\delta U/(2U)\}$ が $N$ に比例しないことと，同じ事実の2つの見え方である。
\end{rikai}
